\documentclass[11pt,a4paper]{article}

\usepackage[utf8]{inputenc}
\usepackage[T1]{fontenc}
\usepackage{lmodern}
\usepackage[a4paper,top=2.4cm,bottom=2.4cm,left=2.6cm,right=2.6cm]{geometry}
\usepackage{parskip}
\usepackage{enumitem}
\usepackage{xcolor}
\usepackage{amsmath}
\usepackage{booktabs}
\usepackage{fancyhdr}
\usepackage{titlesec}

\definecolor{polito}{RGB}{0,60,113}
\definecolor{rulegray}{RGB}{200,200,200}

\titleformat{\section}{\normalfont\large\bfseries\color{polito}}{\thesection.}{0.5em}{}
\titlespacing*{\section}{0pt}{1.5\baselineskip}{0.45\baselineskip}

\pagestyle{fancy}
\fancyhf{}
\renewcommand{\headrulewidth}{0.4pt}
\fancyhead[L]{\footnotesize Thesis summary --- Dylan Magliano (s337915)}
\fancyfoot[C]{\footnotesize \thepage}

\begin{document}

\begin{center}
  {\Large\bfseries\color{polito} POLITECNICO DI TORINO}\\[0.35em]
  {\normalsize Master's Degree in Computer Engineering}\\
  {\small\itshape Artificial Intelligence and Data Analytics track}\\[1.1em]
  {\large\bfseries From pixels to financial portfolios: generating\\
   realistic market scenarios with variational autoencoders}\\[0.7em]
  {\normalsize\bfseries Thesis summary}\\[0.5em]
  {\small Candidate: Dylan Magliano (s337915) \quad---\quad
   Supervisors: prof. Flavio Giobergia, prof. Sergio Caprioli}\\
  {\small Thesis carried out in collaboration with Intesa Sanpaolo
   \quad---\quad Academic Year 2025--2026}
\end{center}

\vspace{0.5em}
\noindent\textcolor{rulegray}{\rule{\textwidth}{0.8pt}}

\section{Context and research problem}

The correlation matrix of asset returns underpins much of quantitative
portfolio analysis, from Markowitz onwards. Estimating it from historical
series alone, however, involves two distinct and structural limitations.

The first is \emph{statistical noise}. When the number of assets $N$ is of the
same order of magnitude as the length $T$ of the series used for estimation,
the empirical matrix becomes unstable and its spectrum fills with spurious
eigenvalues, a phenomenon systematically studied by Random Matrix Theory.

The second, and deeper, is \emph{sampling limitation}: the past offers a single
trajectory among all those that would have been possible. Market
configurations never observed in the specific available sample, yet
plausible in statistical terms, remain inaccessible to any purely historical
analysis.

Hence the research question of this thesis: is it possible to build a deep
generative model that, having learnt how correlations organise themselves in
real data, produces new synthetic correlation matrices --- mathematically
valid and structurally indistinguishable from real ones --- with which to
populate a scenario space richer than the one history alone can provide? The
closest precedent in the literature is CorrGAN (Marti, 2020), which addresses
the same problem with adversarial networks; the work of Caprioli et al.\
(2024) instead uses a Variational Autoencoder on historical matrices with
resampling in the latent space, an approach methodologically akin to the one
adopted here.

\section{Data and pre-processing}

The empirical foundation is a dataset of adjusted closing prices of the
S\&P 500 index over the twenty-year period 2000--2020. To guarantee continuity
of the series and avoid sampling distortions, the universe was filtered by
requiring uninterrupted quotations over the entire horizon: $N = 362$ stocks
survive. The GICS sector classification, required by the hierarchical
structure analyses, was not available from a structured provider and was
reconstructed through automatic classification by a large language model
(360 out of 362 stocks classified).

Daily log-returns are computed from the prices and, from these, Pearson
correlation matrices are estimated with a rolling window procedure. The window
width is set to $T_w = 724$ days, exactly twice the number of assets: this
choice fixes the spectrum ratio $q = N/T_w = 0.5$ by construction, making the
noise level homogeneous over time and analytically governable through the
Marchenko--Pastur bounds, and it guarantees that every matrix is full rank and
therefore positive definite. With a stride of 10 days, 461 matrices are
obtained.

A structural validity filter based on the empirical Perron--Frobenius property
then acts on this raw set, requiring the first eigenvector (the \emph{market
mode}) to have components all of the same sign. The first 49 windows fail it;
detailed analysis shows that these are all and only the windows containing
observations preceding mid-October 2001, in which the only gold-mining stock
in the universe (NEM) is systematically anticorrelated with the rest of the
market in its typical safe-haven role. This is not a measurement artefact but
the signature of an exceptional regime, and excluding it explicitly delimits
the model's domain of validity: \textbf{412 valid matrices} remain, then
randomly shuffled and split 288 / 82 / 42 across training, validation and test
sets. Random shuffling, in place of chronological partitioning, neutralises
the macroeconomic regime \emph{distribution shift} which, in preliminary
tests, prevented the deep architectures from generalising.

\subsection*{The positive semidefiniteness constraint and the Cholesky parametrisation}

A decisive methodological step concerns the form of the input. A correlation
matrix must be positive semidefinite (PSD): this is a global constraint on the
spectrum, and cannot be imposed by acting on individual coefficients
separately. In the first version of the pipeline the models directly
reconstructed the lower triangle of the correlation matrix; spectral checks
revealed the appearance of negative eigenvalues, and a single negative
eigenvalue --- however negligible in magnitude --- is enough to make the
Cholesky decomposition fail, and with it the entire scenario simulation.

The final pipeline solves the problem at its root by changing the object being
learnt: each matrix is decomposed as $\mathbf{C} = \mathbf{L}\mathbf{L}^T$ and
what is flattened and fed to the models is the lower triangular factor
$\mathbf{L}$, not the matrix. Whatever vector the decoder produces, the matrix
$\hat{\mathbf{L}}\hat{\mathbf{L}}^T$ recomposed from it is a Gram matrix and
therefore PSD by construction, independently of the quality of learning; a
final renormalisation to unit diagonal brings the result back into the
correlation domain without altering the signs of the eigenvalues. The
canonical input is therefore a vector of $D = N(N{+}1)/2 = 65\,703$ elements.

\section{Architectures and experimental protocol}

The comparison follows an incremental path across four models, all trained on
the same Cholesky parametrisation and evaluated at equal latent dimension $K$:

\begin{enumerate}[leftmargin=1.4em,itemsep=0.25em,topsep=0.35em]
  \item \textbf{Principal Component Analysis}, a linear benchmark that sets the
        upper limit of linear compression;
  \item \textbf{Linear Autoencoder}, a single dense layer with neither bias nor
        activations, built to replicate that subspace and empirically verify
        the Bourlard and Kamp theorem;
  \item \textbf{Deep Autoencoder}, a funnel design
        $2048 \to 1024 \to 512 \to K$ with LeakyReLU activations, introducing
        non-linearity;
  \item \textbf{Variational Autoencoder}, which shares the funnel encoder of
        the previous model but replaces the deterministic bottleneck with a
        probabilistic inference module ($\boldsymbol{\mu}$,
        $\log\boldsymbol{\sigma}^2$, reparameterization trick).
\end{enumerate}

The VAE reconstruction loss is not the usual mean squared error but the
\textbf{Jeffreys Divergence} between the Gaussian distributions associated
with the original and the reconstructed matrix. The motivation is precise: MSE
treats coefficients as independent quantities and ends up dominated by the
larger correlations, effectively blind to the smaller eigenvalues of the
spectrum --- precisely those that identify minimum-variance portfolios. The
inverse matrices in the Jeffreys formulation reverse this hierarchy and weigh
the error in relative terms along every direction of the spectrum. Because the
reconstruction term is itself a divergence, it is homogeneous with the
Kullback--Leibler regularisation term and the weight $\beta$ remains fixed at
1: the formulation used is therefore the canonical ELBO.

Evaluation combines two families of metrics: spatial errors (MSE, MAE,
Frobenius norm) and topological metrics built on the \emph{Minimum Spanning
Tree} (percentage of shared edges, Jaccard index, overlap of the nodes with
highest betweenness centrality, differences in average path length). The
latter verify that a reconstruction accurate ``on average'' has not nonetheless
altered the hierarchy of links between stocks.

\section{Results on historical reconstruction}

The measurements empirically confirm the Bourlard and Kamp theorem: PCA and the
Linear Autoencoder converge, at equal $K$, on the same error limit. The Deep
Autoencoder clearly surpasses that threshold at every $K$ and shows no sign of
overfitting: at $K = 20$ its test MSE ($7.19 \times 10^{-5}$) is almost four
times lower than PCA's ($2.73 \times 10^{-4}$). The advantage of non-linearity
emerges as one of \emph{representational efficiency}: the AE already reaches
with $K = 10$ the same quality that PCA attains only with $K = 100$.

The Variational Autoencoder sits in an intermediate position. At $K = 20$ it
achieves an MSE of $1.55 \times 10^{-4}$, \textbf{lower} than both PCA and the
Linear Autoencoder, while remaining higher than the deterministic Deep
Autoencoder. MAE and Frobenius norm return exactly the same ranking, which
indicates that the ordering is not an artefact of MSE's sensitivity to outliers.
The topological picture is even sharper and consistent: on MST edge overlap
the VAE recovers two thirds of the distance between the linear models and the
Deep Autoencoder (88.4\% against PCA's 83.0\% and the AE's 91.1\%). The cost
that variational regularisation imposes on reconstruction fidelity is
therefore contained, and smaller than a first intuition would suggest.

Three runs at $K \in \{10, 20, 30\}$, identical in every other hyperparameter,
reveal a peculiar behaviour of the VAE: its reconstruction error \emph{worsens}
slightly as $K$ grows, the opposite of the other three models. Analysis of the
latent space activity spectrum suggests an explanation: in all three configurations only
about six dimensions are effectively active, while the remaining ones undergo
\emph{posterior collapse}. The effective dimensionality of the correlation
manifold therefore appears not to be a hyperparameter to be tuned but rather a property of the
data that the model discovers on its own; adding nominal dimensions does not
increase the capacity actually used and merely widens the space over which
conformity to the prior must be satisfied. The choice of $K = 20$ for the final
model --- the elbow of the reconstruction curve --- is thus robust.

\section{Stochastic scenario generation and validation}

It is on the generative side that the VAE shows its distinctive value. The
deterministic Deep Autoencoder, however superior in pure reconstruction, has a
fragmented latent space devoid of meaning outside the observed points:
sampling a random vector from it produces output with no economic or financial
coherence. Kullback--Leibler regularisation, combined with sampling noise,
instead shapes a continuous and dense latent space in which meaningful
sampling is possible. The most effective synthesis of the results is this
complementarity: \emph{the AE to remember, the VAE to imagine}.

The generation procedure actually adopted, however, is not a simple sampling
from the prior $\mathcal{N}(\mathbf{0},\mathbf{I})$, which would produce some
arbitrary plausible scenario, disconnected from the present. The objective is
more specific: to build a distribution of realistic scenarios of the
correlation structure \emph{at a one-month horizon starting from the market
state observed today}, to be assessed from a risk perspective. The technique
used is a \textbf{historical block bootstrap in the latent space}, articulated
in four steps: the encoder reconstructs the historical latent trajectory day by
day; its daily increments are computed; a block of 21 consecutive increments
(one working month) is randomly sampled and its cumulative jump computed; the
jump is added to the current state and decoded, yielding a PSD matrix by
construction. Repeating this a thousand times yields a distribution of possible
evolutions, anchored to the present yet diversified by block-sampled historical
variability.

The one thousand scenarios so produced were validated against the five
stylized facts of financial correlation matrices, each translated into a
scalar indicator comparable with today's matrix and with the range observed
over the 412 real matrices:

\begin{enumerate}[leftmargin=1.4em,itemsep=0.2em,topsep=0.35em]
  \item \textbf{Pairwise distribution skewed towards positive values}: mean
        between 0.4609 and 0.5352 across all scenarios, against 0.4512 for
        today's matrix.
  \item \textbf{Spectrum and Marchenko--Pastur}: market eigenvalue $\lambda_1$
        in the range $[172.4,\ 197.2]$ against today's 169.1, roughly two
        orders of magnitude beyond the upper edge of the bulk, accompanied by
        7--9 sector modes.
  \item \textbf{Perron--Frobenius property}: the minimum component of the
        market mode is always positive, with a larger margin from the violation
        threshold than today's matrix itself.
  \item \textbf{Hierarchical sector structure}: the ratio between average
        intra-sector and inter-sector correlation is always greater than 1 and
        within the historical range; the dendrograms visually reproduce the
        same sector colour blocks as the real matrix.
  \item \textbf{Scale-free MST topology}: the power-law exponent lies between
        2.053 and 2.306, inside the range $2 < \gamma < 3$ expected from the
        literature.
\end{enumerate}

The outcome is that \textbf{all 1000 scenarios satisfy the five criteria
simultaneously}, and all matrices are positive definite, with positive semidefiniteness guaranteed by
construction by the Cholesky parametrisation. Detailed graphical inspection of
a single scenario corroborates the picture on those objects no scalar indicator can
summarise: cosine similarity with respect to today's matrix is 0.9999 on the
market mode and 0.9893 and 0.9883 on the two subsequent sector modes.

\section{Limitations and future work}

The work explicitly declares its methodological limitations. The aggregate
verification reduces each stylized fact to a scalar indicator, while full
visual inspection remains limited to a single representative scenario; the one
thousand scenarios moreover derive from heavily overlapping 21-day blocks, with
an effective sample size below one thousand, and the entire verification is
anchored to a single starting market state. Random shuffling, while addressing
distribution shift, does not remove the informational overlap between adjacent
windows (up to 98.6\% of shared observations), which represents a limited but
real risk of optimism in the reported generalisation metrics. Finally, the
model comparison does not cover all latent dimensions uniformly for the VAE,
and the GICS sector classification was not validated against official sources.

The natural continuation is to apply the generated scenario distribution to
portfolio risk management, comparing the resulting picture with the one
obtainable from historical data alone. Next come the strengthening of
statistical validation --- repeating the verification from a plurality of
initial states distributed across the twenty-year horizon and extending it from
scalar indicators to formal distributional goodness-of-fit tests --- a direct
comparison with alternative generative paradigms (CorrGAN, and more generally
diffusion models), and the extension of the pipeline to other asset classes and
to international equity universes.

\section{Conclusion}

The thesis has designed, implemented and empirically evaluated a complete framework for the
compression, reconstruction and stochastic generation of realistic financial
correlation matrices. The most significant result is that the Cholesky
parametrisation structurally resolves the positive semidefiniteness constraint,
and that on this basis the Variational Autoencoder pays a contained cost in
reconstruction fidelity --- indeed remaining superior to the linear benchmarks
--- in exchange for the property that makes stochastic generation possible.
The simultaneous conformity of the one thousand generated scenarios to the five
stylized facts supports the proposed architecture as a self-contained tool for
producing mathematically coherent synthetic financial ecosystems.

\end{document}
